% Lecture visuelle directe avec la vision du modèle ; aucun OCR.
% Source : PDF original, vue 004, côté gauche, page manuscrite 6.
% Statut : transcrit ; relecture indépendante à effectuer.
\chapter*{Plan de tout l'ouvrage}
\addcontentsline{toc}{chapter}{Plan de tout l'ouvrage}
\markboth{Plan de tout l'ouvrage}{Plan de tout l'ouvrage}

L'homme créé à l'image de Dieu porte en lui-même, Messieurs et chers amis, quelque chose qui lui fait comprendre que l'immortalité doit être un jour son partage... C'est à cause de cela que son imagination, toujours active, voudroit embrasser à la fois le passé et l'avenir. Nous voudrions qu'il nous fût donné de connoître nos destinées futures ; nous voudrions encore pouvoir retrouver à travers les siècles la trace du chemin qui nous a conduits au point où nous sommes parvenus.

De là cette disposition naturelle non seulement aux grandes nations, dont les annales se lient à l'histoire de l'univers, mais encore aux plus minces provinces, aux plus petites tribus, aux villes, aux bourgades, et même à de simples aggrégations, à rechercher tout ce qui concerne leur origine et leur histoire. Si ce sentiment qui se trouve au fond de tous les cœurs mérite qu'on y fasse quelque attention lorsqu'il n'est pour toute question que d'un établissement purement humain, à plus forte raison est-il juste et raisonnable de le satisfaire lorsqu'il a pour but de faire connoître les moyens que la Providence a employés soit pour donner naissance à une société destinée à concourir au salut et le bien des personnes, soit pour la soutenir et la conserver au milieu des dangers qui ont pu menacer son existence. Il entre donc dans l'ordre que la Congrégation de Lyon, si évidemment l'œuvre de Dieu, désire de posséder son histoire afin de pouvoir transmettre aux générations futures le souvenir des grâces dont elle a
% Source : PDF original, vue 004, côté droit, page manuscrite 7.
été l'objet. Il est temps d'aviser aux moyens de satisfaire promptement ce pieux désir : il y aura bientôt 40 ans que la Congrégation a été fondée ; la mort a déjà moissonné plusieurs de ses premiers membres. On peut entrevoir le moment assez prochain où les témoins de ses premières années auront tous disparu de la scène de ce monde ; et cependant tous les documens écrits de ces premiers temps ayant été anéantis, comme nous aurons à le raconter, ce n'est qu'en interrogeant les souvenirs des anciens qu'on peut parvenir à conserver cette si précieuse histoire. C'est cette pensée qui m'a déterminé, pendant que je le puis encore, à coopérer, autant que cela dépendra de moi, à l'établissement de ces annales.

Voici le plan que je me suis formé : je vais d'abord rassembler dans un précis historique les divers événements, soit publics, soit particuliers, qui m'ont paru avoir donné lieu à la fondation de la Congrégation. Je ferai entrer dans ce précis même les circonstances de détail qui ont eu quelque part à cette fondation. Cela donnera à la Congrégation une connaissance exacte des moyens dont la Providence a voulu se servir pour préparer notre union. Si je ne puis pas dans ce précis vous tout faire connoître complètement, du moins j'espère vous transmettre à ce sujet tout ce que j'ai vu de mes propres yeux ; je n'omettrai rien de ce qui sera parvenu à ma connoissance. J'espère pouvoir ensuite vous donner une histoire assez exacte de la Congrégation des jeunes gens depuis sa fondation jusqu'à l'année 1817, époque de la fondation de celle des hommes. Mais arrivé à cette nouvelle période l'histoire se complique, et trop de documens me manquent, ne pouvant les trouver que dans mes souvenirs, pour qu'il me
% Source : PDF original, vue 005, côté gauche, page manuscrite 8.
fût possible de vous offrir une histoire complète de cette dernière partie. Du reste ce sont heureusement les temps les plus modernes, dont je pourrai du moins vous rendre compte, et avec le secours des procès-verbaux de l'époque qui subsistent et les notes qu'ont conservées MM. les Préfets, il sera facile à une plume plus exercée que la mienne de terminer cette histoire.

Je vais donc me borner quant à présent à vous donner sur ces derniers temps de simples mémoires historiques, qui feront suite à l'histoire des premiers temps. Autant que mes souvenirs me le permettront, je m'efforcerai de vous rappeler surtout les différentes bonnes œuvres auxquelles la Congrégation a pris part, et toutes les circonstances de son histoire qui offrent quelque intérêt et qui n'étant pas de nature à laisser des traces dans les procès-verbaux. Ce seront des documens préparés pour l'histoire future.

Ainsi : 1\up{o} précis historique des faits antérieurs à la fondation. 2\up{o} histoire des quinze premières années. 3\up{o} mémoires historiques sur les temps qui ont suivi. Voilà tout le plan que je me propose de suivre avec la grâce de Dieu et sous les auspices de Marie. Entrons donc en matière.
